\section{Experiments}
\label{sec:experiments}

The experiments test one sealed score in two distinct roles. RQ1 asks whether
it prospectively orders later damage. RQ2 tests, on each setting's declared
fidelity support, whether its forward-only coordinate recovers the ordering
induced by exact Loss Susceptibility. RQ3 asks whether
the same score improves a real allocation decision under a fixed repair
budget. Tables~\ref{tab:pred-value}--\ref{tab:prot-contract} report the three
primary claims; Table~\ref{tab:robustness} reports coverage and failure
boundaries.

\subsection{Experimental Design and Evidence Rules}
\label{sec:benchmark-protocol}

\paragraph{Sealed data roles.}
Candidates are split by semantic group into discovery and audit folds.
Three disjoint, target-request-disjoint development folds separately fix the
operating point $(R,\eta)$ for the Loss Susceptibility estimator
and fidelity fraction $q_{\mathrm{fid}}$, the shared mixture weight
$\widehat\alpha_{\mathrm{pred}}$, and the repair protocol. The latter includes
$K_p$, $B_{\mathrm{tok}}$, stopping rules, and guards. On a target
request, discovery candidates may fit rank maps and supply repair examples;
audit candidates never enter repair. No target damage or audit outcome can
change the score, pool, parent checkpoint, or returned repair checkpoint.

\paragraph{Settings and parents.}
The campaign spans TOFU \cite{maini2024tofu}, WMDP-bio/MMLU
\cite{li2024wmdp,hendrycks2021measuring}, MUSE-News \cite{shi2024muse}, and
RWKU \cite{jin2024rwku}; MUSE-Books and PISTOL \cite{qiu2024pistol} probe
boundary conditions. Models include Qwen2.5-1.5B, Qwen2.5-7B, Qwen2.5-14B,
and Llama-3.1-8B. The primary parent roster is GradDiff, NPO, SimNPO, GRU,
RMU, RepNoise, and Circuit Breakers, with GA and IdkDPO retained as stress
controls. A campaign manifest freezes every
model--dataset--request--parent cell and seed.

\paragraph{Prediction references.}
The deployed score $S_{\widehat\alpha_{\mathrm{pred}}}$ is compared with its
two endpoints: $S_0$ uses the rank-normalized Loss Susceptibility
estimate only, and $S_1$ uses Representation Proximity only. Additional references include exact
Loss Susceptibility, checkpoint-local candidate--forget-gradient
alignment computed at $\theta_0$, initial NLL, answer length, lexical and
sentence-embedding proximity, request-shuffled
Representation Proximity, a one-direction loss response, and repeated random
rankings. The strongest simple control is chosen on development requests and
frozen before target evaluation.

\paragraph{Protection comparators.}
The two-signal selector, the two endpoint selectors, and repeated random
selection use exactly the same repair harness. The unrepaired entry checkpoint
measures net repair benefit. The $q_G^\star$-substituted mixture selector,
uniform rehearsal, and a
realized-discovery-damage oracle are diagnostic arms that separate estimator
error, repair capacity, and selector quality; they do not replace the
confirmatory comparisons.
Retain-aware mitigation methods that reweight or curate the forget data,
constrain the parent objective, or rectify its gradients intervene on the
parent trajectory itself; running them as arms would confound selector quality
with the operator and objective differences this harness deliberately holds
fixed. They are therefore contrasted in the related-work discussion rather
than included as confirmatory comparators.

\paragraph{Outcomes and aggregation.}
The primary retained-behavior outcome is candidate-level $\Delta$NLL relative
to $\theta_0$. We report request-level mean damage and
$\mathrm{CVaR}_{.95}$, the mean of the worst 5\% of candidate damage, together
with each dataset's native retained metric and forgetting/utility slack.
Prediction associations are Spearman rank correlations computed within each
request--seed--parent cell. Seeds are averaged within requests and requests
receive equal weight. A paired hierarchical bootstrap resamples requests,
trajectory seeds, and semantic groups.

\paragraph{How the three claims are decided.}
RQ1 requires a positive prospective rank association, paired improvement over
both signal endpoints and the strongest simple control, and above-chance
recovery of the harmful tail on sufficient support. For compact table
notation, $g_G$ is the gain over $S_0$, $g_H$ is the gain over $S_1$, and
$g_{\mathrm{ctl}}$ is the gain over the frozen simple control.
Table~\ref{tab:pred-value} additionally reports the frozen mixture weight
$\widehat\alpha_{\mathrm{pred}}$ of every cell; an unresolved fallback weight
is claim-ineligible and marked~$\dagger$.

RQ2 is a setting-level numerical-estimator claim. It uses rank agreement
$f_\rho$ between estimated and exact Loss Susceptibility,
$\widehat q_G$ and $q_G^\star$, and their Top-$K_{\mathrm{fid}}$ overlap $f_K$
on the setting's declared fidelity support. The one-sided bounds for $f_\rho$
and $f_K$ resample that support. Each target profile separately uses
split-bank rank agreement, perturbation survival, and integrity coverage as
outcome-blind eligibility checks; these checks are not an exact-rank
certificate for its candidate support. Time and memory are also reported.
Predictive usefulness belongs to RQ1 rather than RQ2.

RQ3 compares the two-signal repair arm with each single-signal arm, repeated
random selection, and no repair. It requires lower held-out mean and tail
damage, non-inferiority on the native retained metric, and complete common support
inside the same forgetting/utility feasibility region. Eligibility and
statistical success are reported separately. Appendix~\ref{app:outcomes}
specifies thresholds, missingness, multiplicity, and the complete decision
rules.

% Stable include path for the evidence renderer's generated core tables.
% Generated by experiments/paper/build_evidence.py; do not edit by hand.
% Incomplete evidence remains an explicit \tblph placeholder.
\begin{table*}[!t]
\caption{\textbf{Prospective damage prediction (RQ1, rank condition):
the sealed mixture against its own ingredients.} Scores are sealed
before parent runs. Point correlations for both endpoints, the sealed
mixture with its one-sided 95\% lower bound, and paired gains over
both endpoints and the frozen simple control are shown.
\texttt{E/P} denotes eligible/pass.}
\label{tab:pred-value}
\centering
\small
\setlength{\tabcolsep}{5pt}
\resizebox{\textwidth}{!}{%
\begin{tabular}{@{}lcccccccc@{}}
\toprule
\multicolumn{9}{@{}l}{\textbf{Prospective prediction loss-shake validation (TOFU, Qwen2.5-7B)}} \\
Parent & $\widehat\alpha_{\mathrm{pred}}$ & $S_0$ & $S_1$ (prox.) & $S_{\widehat\alpha_{\mathrm{pred}}}$ (ours) [LB] & $g_G$ [LB] & $g_H$ [LB] & $g_{\rm ctl}$ [LB] & Rank E/P \\
\midrule
\multicolumn{9}{@{}l}{\emph{Output-readout parents}} \\
GradDiff & \tblph & \tblph & \tblph & \tblph & \tblph & \tblph & \tblph & n/-- \\
NPO & \tblph & \tblph & \tblph & \tblph & \tblph & \tblph & \tblph & n/-- \\
SimNPO & \tblph & \tblph & \tblph & \tblph & \tblph & \tblph & \tblph & n/-- \\
GRU & \tblph & \tblph & \tblph & \tblph & \tblph & \tblph & \tblph & n/-- \\
\addlinespace[2pt]
\multicolumn{9}{@{}l}{\emph{Representation-readout parents}} \\
RMU & \tblph & \tblph & \tblph & \tblph & \tblph & \tblph & \tblph & n/-- \\
RepNoise & \tblph & \tblph & \tblph & \tblph & \tblph & \tblph & \tblph & n/-- \\
CB & \tblph & \tblph & \tblph & \tblph & \tblph & \tblph & \tblph & n/-- \\
\bottomrule
\end{tabular}}
\end{table*}

\begin{table*}[t]
\caption{\textbf{Loss Susceptibility estimation fidelity against
its exact target (RQ2).} Panel A reports the registered setting-level
certificate. Panel B substitutes exact gradient energy inside the
frozen profile without refitting and is diagnostic only.}
\label{tab:loss-susceptibility-fidelity}
\centering
\footnotesize
\setlength{\tabcolsep}{5pt}
\resizebox{\textwidth}{!}{%
\begin{tabular}{@{}llccccccc@{}}
\toprule
\multicolumn{9}{@{}l}{\textbf{A. Estimator versus direct-gradient reference on declared setting-level fidelity support}} \\
Estimator & Access & $f_\rho$ [LB$-\tau_\rho$] & $f_K$ [LB$-\tau_K$] & Split-bank $\rho$ & Survival & Time (s) & Mem (GB) & RQ2 E/P \\
\midrule
Direct-gradient reference & reverse, per-cand. & 1.00 [--] & 1.00 [--] & -- & -- & \tblph & \tblph & -- \\
Loss Susceptibility, $2R$ forward sweeps (\textbf{ours}) & forward-only & \tblph & \tblph & \tblph & \tblph & \tblph & \tblph & n/-- \\
\bottomrule
\end{tabular}}

\vspace{6pt}
\begin{tabular}{@{}lcc@{}}
\toprule
\multicolumn{3}{@{}l}{\textbf{B. $q_G^\star$-substituted mixture
inside the frozen profile (diagnostic only)}} \\
\midrule
Parent & $\Delta\rho_{\mathrm{swap}}$ [95\% CI] & $g_H$ [LB] (Tab.~\ref{tab:pred-value}) \\
\midrule
\multicolumn{3}{@{}l}{\emph{Output-readout parents}} \\
GradDiff & \tblph & \tblph \\
NPO & \tblph & \tblph \\
SimNPO & \tblph & \tblph \\
GRU & \tblph & \tblph \\
\addlinespace[2pt]
\multicolumn{3}{@{}l}{\emph{Representation-readout parents}} \\
RMU & \tblph & \tblph \\
RepNoise & \tblph & \tblph \\
CB & \tblph & \tblph \\
\bottomrule
\end{tabular}
\end{table*}

\begin{table*}[t]
\caption{\textbf{Harmful-tail recovery (RQ1, tail condition).}
The harmful tail contains the largest positive realized damages.
Endpoint and sealed-mixture recall are compared with exact finite-sample
chance; lift carries its one-sided 95\% lower bound.}
\label{tab:tail-recovery}
\centering
\small
\setlength{\tabcolsep}{5pt}
\begin{tabular}{@{}lccccccc@{}}
\toprule
Parent & Tail-eligible $n/N$ & Chance $q_{\mathrm{tail}}^{\mathrm{eff}}$ & $S_0$ & $S_1$ (prox.) & $S_{\widehat\alpha_{\mathrm{pred}}}$ (ours) & $L_{\mathrm{tail}}$ [LB] & RQ1 E/P \\
\midrule
\multicolumn{8}{@{}l}{\emph{Output-readout parents}} \\
GradDiff & \tblph & \tblph & \tblph & \tblph & \tblph & \tblph & n/-- \\
NPO & \tblph & \tblph & \tblph & \tblph & \tblph & \tblph & n/-- \\
SimNPO & \tblph & \tblph & \tblph & \tblph & \tblph & \tblph & n/-- \\
GRU & \tblph & \tblph & \tblph & \tblph & \tblph & \tblph & n/-- \\
\addlinespace[2pt]
\multicolumn{8}{@{}l}{\emph{Representation-readout parents}} \\
RMU & \tblph & \tblph & \tblph & \tblph & \tblph & \tblph & n/-- \\
RepNoise & \tblph & \tblph & \tblph & \tblph & \tblph & \tblph & n/-- \\
CB & \tblph & \tblph & \tblph & \tblph & \tblph & \tblph & n/-- \\
\bottomrule
\end{tabular}
\end{table*}

\begin{table*}[t]
\caption{\textbf{Fixed-budget protection (RQ3, effect condition):
allocation is the only free variable.} Damage cells report
mean;\,$\mathrm{CVaR}_{.95}$ on common held-out support.}
\label{tab:prot-effect}
\centering
\small
\setlength{\tabcolsep}{5pt}
\begin{tabular}{@{}lccccccc@{}}
\toprule
Parent & No repair & Random & $S_0$ & $S_1$ (prox.) & $S_{\widehat\alpha_{\mathrm{pred}}}$ (ours) & $\max\widehat\Delta_{a,k}$ / $\max U^{95}_{a,k}$ & Effect E/P \\
\midrule
\multicolumn{8}{@{}l}{\emph{Output-readout parents}} \\
GradDiff & \tblph & \tblph & \tblph & \tblph & \tblph & \tblph & n/-- \\
NPO & \tblph & \tblph & \tblph & \tblph & \tblph & \tblph & n/-- \\
SimNPO & \tblph & \tblph & \tblph & \tblph & \tblph & \tblph & n/-- \\
GRU & \tblph & \tblph & \tblph & \tblph & \tblph & \tblph & n/-- \\
\addlinespace[2pt]
\multicolumn{8}{@{}l}{\emph{Representation-readout parents}} \\
RMU & \tblph & \tblph & \tblph & \tblph & \tblph & \tblph & n/-- \\
RepNoise & \tblph & \tblph & \tblph & \tblph & \tblph & \tblph & n/-- \\
CB & \tblph & \tblph & \tblph & \tblph & \tblph & \tblph & n/-- \\
\bottomrule
\end{tabular}
\end{table*}

\begin{table*}[t]
\caption{\textbf{Fixed-budget protection (RQ3, contract condition):
the comparison is constraint-matched.} Native non-inferiority,
forgetting and utility slack, common-support feasibility, and guarded
repair updates are reported separately from the damage effect.}
\label{tab:prot-contract}
\centering
\small
\setlength{\tabcolsep}{5pt}
\begin{tabular}{@{}lcccccc@{}}
\toprule
Parent & $\min\widehat h_a$ / $\min L^{95}_a$ & Forget slack & Utility slack & Feasible arms & Updates/rollback & RQ3 E/P \\
\midrule
\multicolumn{7}{@{}l}{\emph{Output-readout parents}} \\
GradDiff & \tblph & \tblph & \tblph & \tblph & \tblph & n/-- \\
NPO & \tblph & \tblph & \tblph & \tblph & \tblph & n/-- \\
SimNPO & \tblph & \tblph & \tblph & \tblph & \tblph & n/-- \\
GRU & \tblph & \tblph & \tblph & \tblph & \tblph & n/-- \\
\addlinespace[2pt]
\multicolumn{7}{@{}l}{\emph{Representation-readout parents}} \\
RMU & \tblph & \tblph & \tblph & \tblph & \tblph & n/-- \\
RepNoise & \tblph & \tblph & \tblph & \tblph & \tblph & n/-- \\
CB & \tblph & \tblph & \tblph & \tblph & \tblph & n/-- \\
\bottomrule
\end{tabular}
\end{table*}


% Stable include path for the evidence renderer's generated robustness tables.
% Generated-table shell; experiments/paper/build_evidence.py owns this path.
\begin{table*}[!t]
  \caption{\textbf{Breadth of the claim across predeclared settings.} Each
  row is one setting on a declared breadth axis (held-out requests, datasets,
  model scale/family); its seven parents are fixed in advance, and
  \emph{Plan/done} keeps unfinished rows in the denominator. The RQ1 and RQ3
  columns report how many of the seven parents were eligible and passed
  (\texttt{E/P}); RQ2 reports the setting's single registered fidelity
  certificate, evaluated on its declared setting-level support and shared by
  all parents in that setting.
  Eligibility failures are diagnosed in
  Table~\ref{tab:robustness-funnel}, and \textit{n.a.}\ marks a setting with
  no registered fidelity certificate. \emph{Chain} is the licensing rule: a
  setting supports the full claim only when its setting-level certificate
  passes RQ2, every claim-bearing target profile passes the separate
  split-bank, perturbation-survival, and integrity eligibility checks, and at
  least one output-readout and one representation-readout parent are eligible
  and pass both RQ1 and RQ3 after the predeclared within-readout correction.
  The certificate does not exact-rank certify each target candidate support.
  The paper-level breadth claim
  additionally requires a licensed chain in primary TOFU, in one non-primary
  model scale or family, and in at least two of WMDP-bio, MUSE-News, and
  RWKU; stress settings (marked) can never rescue a failure. We do not pool
  parents into a universal effect.}
  \label{tab:robustness}
  \centering
  \small
  \setlength{\tabcolsep}{5pt}
  \begin{tabular}{@{}llccccc@{}}
    \toprule
    & & & \multicolumn{3}{c}{Eligible/passing counts} & \\
    \cmidrule(lr){4-6}
    Axis & Setting
      & Plan/done
      & RQ1 parents
      & RQ2 cert.
      & RQ3 parents
      & Chain \\
    \midrule
    Request & held-out TOFU requests (Qwen2.5-7B) & 7/0 & 0/0 & 0/0 & 0/0 & -- \\
    \addlinespace[2pt]
    Dataset & WMDP-bio/MMLU        & 7/0 & 0/0 & 0/0 & 0/0 & -- \\
            & MUSE-News            & 7/0 & 0/0 & \textit{n.a.} & 0/0 & -- \\
            & RWKU                 & 7/0 & 0/0 & 0/0 & 0/0 & -- \\
            & MUSE-Books (stress)  & 7/0 & 0/0 & \textit{n.a.} & 0/0 & -- \\
            & PISTOL (stress)      & 7/0 & 0/0 & \textit{n.a.} & 0/0 & -- \\
    \addlinespace[2pt]
    Model   & Qwen2.5-1.5B (boundary) & 7/0 & 0/0 & \textit{n.a.} & 0/0 & -- \\
            & Qwen2.5-14B          & 7/0 & 0/0 & 0/0 & 0/0 & -- \\
            & Llama-3.1-8B         & 7/0 & 0/0 & 0/0 & 0/0 & -- \\
    \bottomrule
  \end{tabular}
\end{table*}

\begin{table*}[!t]
  \caption{\textbf{Predeclared evidence funnel and failure boundaries.} For
  each setting, columns follow the stages a setting--parent row must survive
  before any claim can pass, so a drop between adjacent columns localizes the
  failure: sealed profiles valid (of planned), trajectories reaching the
  common forgetting gate (of planned), audit cells with complete common
  support for prediction, tail-eligible cells among them, and rows with all
  five repair arms feasible (of reached rows with a valid profile). No stage
  removes a row from the declared denominator. \emph{Worst bounds} are the
  least-favorable one-sided bounds among attempted rows --- the minimum RQ1
  and RQ2 lower bounds and the maximum individual RQ3 damage upper bound --- reported as
  descriptive extrema of how close the setting came to the boundary, not as
  claims. \emph{Failure modes} name the dominant predeclared reasons
  (invalid profile, non-reach, no common support, tail miss, infeasible arm,
  fallback weight, \ldots); \emph{not attempted} marks settings whose
  campaign has not started.}
  \label{tab:robustness-funnel}
  \centering
  \small
  \setlength{\tabcolsep}{5pt}
  \resizebox{\textwidth}{!}{%
  \begin{tabular}{@{}llccccccl@{}}
    \toprule
    & & \multicolumn{5}{c}{Survival down the sealed pipeline} & & \\
    \cmidrule(lr){3-7}
    Axis & Setting
      & Profiles valid
      & Gate reached
      & Common $n$
      & Tail-elig.\ $n$
      & All-arm feas.
      & Worst RQ1/RQ2/RQ3
      & Failure modes \\
    \midrule
    Request & held-out TOFU requests (Qwen2.5-7B) & -- & -- & -- & -- & -- & -- & not attempted \\
    \addlinespace[2pt]
    Dataset & WMDP-bio/MMLU        & -- & -- & -- & -- & -- & -- & not attempted \\
            & MUSE-News            & -- & -- & -- & -- & -- & -- & not attempted \\
            & RWKU                 & -- & -- & -- & -- & -- & -- & not attempted \\
            & MUSE-Books (stress)  & -- & -- & -- & -- & -- & -- & not attempted \\
            & PISTOL (stress)      & -- & -- & -- & -- & -- & -- & not attempted \\
    \addlinespace[2pt]
    Model   & Qwen2.5-1.5B (boundary) & -- & -- & -- & -- & -- & -- & not attempted \\
            & Qwen2.5-14B          & -- & -- & -- & -- & -- & -- & not attempted \\
            & Llama-3.1-8B         & -- & -- & -- & -- & -- & -- & not attempted \\
    \bottomrule
  \end{tabular}}
\end{table*}



\subsection{RQ1: Does the Sealed Score Predict Later Damage?}
\label{sec:exp-prediction}

\paragraph{Answer.}
{\setlength{\emergencystretch}{3em}
\PredictionHeadline\ \TailHeadline\par}

\paragraph{Candidate-level ordering.}
We compare the score sealed at $\theta_0$ with damage revealed at
$\theta_{t^\dagger}$ on the untouched audit fold. The primary evidence is its
within-request rank correlation and its paired gains over $S_0$, $S_1$, and
the frozen simple control on identical candidate support. Reporting the
mixture weight, the absolute association, and all three gains prevents a
strong endpoint or a trivial checkpoint feature from being hidden by the
combined score.

\paragraph{Harmful-tail recovery.}
We also test whether the most damaged behaviors are found early enough to
prioritize. We predeclare an audit-tail fraction
$q_{\mathrm{tail}}$ and compare the top
$K_{\mathrm{tail}}=\lceil q_{\mathrm{tail}}
|\mathcal C_{\mathrm{audit}}|\rceil$ sealed ranks with the
$K_{\mathrm{tail}}$ largest positive-damage outcomes. We report recall
relative to its exact finite-sample chance level, together with
positive-damage support and tail-eligibility coverage. This prediction-tail
measure is distinct from the $\mathrm{CVaR}_{.95}$ outcome used to evaluate
repair in RQ3.

\paragraph{Evidence and boundaries.}
Tables~\ref{tab:pred-value} and~\ref{tab:tail-recovery} report the primary
prediction quantities by parent. Appendix Table~\ref{tab:prospective-prediction} expands the predictor
roster, horizons, and construction checks. Reversals against an endpoint,
diffuse or unsupported tails, invalid profiles, and gate-non-reaching runs
remain visible rather than being averaged away.

\subsection{RQ2: Is Loss Susceptibility Faithfully Estimated?}
\label{sec:exp-fidelity}

\paragraph{Answer.}
{\setlength{\emergencystretch}{3em}\FidelityHeadline\par}

\paragraph{Ranking and screening fidelity.}
At the operating point frozen on the calibration fold, we compare
$\widehat q_G$ with exact Loss Susceptibility $q_G^\star$ on the
same block and declared setting-level fidelity support. Rank agreement measures
agreement across the full ranking, while Top-$K_{\mathrm{fid}}$ overlap tests whether the
high-susceptibility set relevant to screening is recovered. A predeclared
nested-$R$ frontier
shows how both quantities change with the number of shared directions.

\paragraph{Stability and access cost.}
The fidelity report includes perturbation survival, agreement between
disjoint halves of the direction bank, synchronized time, peak memory, and
valid-profile coverage. The direct-gradient reference streams one candidate
gradient at a time; the Loss Susceptibility estimator uses $2R$
batched pool-forward sweeps and no
candidate-side backward pass during target profiling. The intended claim is
faithful rank recovery on the declared fidelity support together with
backward-free target profiling under separate numerical eligibility checks,
not guaranteed wall-clock superiority or exact-rank certification of every
target support.

\paragraph{What RQ2 does not claim.}
RQ2 supports only the rank-based estimator claim at the frozen operating point
on the declared fidelity support. Passing it shows that $\widehat q_G$
recovers the ordering and high-susceptibility set induced by $q_G^\star$ on
that support. It neither exact-rank certifies a target candidate support nor
establishes magnitude accuracy. Nor does RQ2 establish that Loss Susceptibility
predicts unlearning damage, which the endpoint and control
comparisons test in RQ1.
Table~\ref{tab:loss-susceptibility-fidelity} reports the decisive fidelity quantities;
Appendix Table~\ref{tab:bwfree} gives the full fidelity--cost frontier and
integrity checks.

\subsection{RQ3: Does the Sealed Score Improve Repair Allocation?}
\label{sec:exp-protection}
\label{sec:exp-alpha}

\paragraph{Answer.}
{\setlength{\emergencystretch}{3em}\ProtectionHeadline\par}

\paragraph{Selector comparison.}
Each confirmatory repair arm starts from the common criterion-reaching parent
checkpoint. They share the repair operator, neutral stream, example order,
optimizer random stream, token budget, guards, and save schedule; only the
$K_p$ selected discovery behaviors change. Panel C of
Figure~\ref{fig:profile-two-uses} summarizes this controlled comparison.

The confirmatory outcomes are mean and $\mathrm{CVaR}_{.95}$ damage on the
untouched audit fold. Lower values are better.
Tables~\ref{tab:prot-effect} and~\ref{tab:prot-contract} report our arm's
held-out damage beside the
strongest alternative selector and the unrepaired checkpoint, together with
the worst point estimates and worst individual one-sided 95\% bounds for paired
damage and native-metric non-inferiority, plus the minimum forgetting/utility
slack. This separates selector quality from
differences in repair compute or checkpoint admissibility.

\paragraph{Net repair benefit.}
Single-signal and random selectors test allocation quality under the same
repair capacity. The no-repair checkpoint answers a
different but necessary question: whether the repair pipeline provides net
benefit at all. A selector cannot pass RQ3 merely by winning among repaired
arms if repair itself worsens held-out behavior or violates the common
contract.

\paragraph{Diagnostics and boundaries.}
The $q_G^\star$-substituted mixture selector, uniform rehearsal, and the
realized-damage oracle remain secondary diagnostics outside RQ3. Appendix
Table~\ref{tab:crossed} reports the predeclared selection-quota and specificity
boundary checks. Infeasible arms, incomplete random-draw
rosters, missing native outcomes, and forgetting or utility regressions remain
failures rather than being dropped.

\subsection{Breadth and Failure Boundaries}
\label{sec:exp-supporting}

\paragraph{Answer.}
{\setlength{\emergencystretch}{3em}\TransferHeadline\par}

Table~\ref{tab:robustness} keeps every predeclared model, benchmark, request,
and parent in view instead of pooling them into one universal effect. It
reports planned, completed, eligible, and passing cells alongside invalid
profiles, parent non-reach, endpoint reversals, diffuse damage, infeasible
repair, and contract regressions. Breadth requires evidence in both output-
and representation-readout parent families and replication across the
predeclared model and benchmark settings. Readout labels describe objectives;
they never choose a signal endpoint or alter a target rule.
